\documentclass[10pt,a4paper]{amsart}
\usepackage{amsmath,amssymb,amsthm,graphicx}
% For commutative diagrams we use Xy-pic
\usepackage[all]{xy}

% We use 2cell for 2-commutative diagrams.
\xyoption{2cell}
% Theorem environments.
%
\theoremstyle{plain}
\newtheorem{theorem}[subsection]{Theorem}
\newtheorem{proposition}[subsection]{Proposition}
\newtheorem{lemma}[subsection]{Lemma}

\theoremstyle{definition}
\newtheorem{definition}[subsection]{Definition}
\newtheorem{example}[subsection]{Example}
\newtheorem{exercise}[subsection]{Exercise}
\newtheorem{situation}[subsection]{Situation}

\theoremstyle{remark}
\newtheorem{remark}[subsection]{Remark}
\newtheorem{remarks}[subsection]{Remarks}

\numberwithin{equation}{subsection}

% Macros
%
\def\lim{\mathop{\mathrm{lim}}\nolimits}
\def\colim{\mathop{\mathrm{colim}}\nolimits}
\def\Spec{\mathop{\mathrm{Spec}}}
\def\Hom{\mathop{\mathrm{Hom}}\nolimits}
\def\Ext{\mathop{\mathrm{Ext}}\nolimits}
\def\SheafHom{\mathop{\mathcal{H}\!\mathit{om}}\nolimits}
\def\SheafExt{\mathop{\mathcal{E}\!\mathit{xt}}\nolimits}
\def\Sch{\mathit{Sch}}
\def\Mor{\mathop{\mathrm{Mor}}\nolimits}
\def\Ob{\mathop{\mathrm{Ob}}\nolimits}
\def\Sh{\mathop{\mathit{Sh}}\nolimits}
\def\NL{\mathop{N\!L}\nolimits}
\def\CH{\mathop{\mathrm{CH}}\nolimits}
\def\proetale{{pro\text{-}\acute{e}tale}}
\def\etale{{\acute{e}tale}}
\def\QCoh{\mathit{QCoh}}
\def\Ker{\mathop{\mathrm{Ker}}}
\def\Im{\mathop{\mathrm{Im}}}
\def\Coker{\mathop{\mathrm{Coker}}}
\def\Coim{\mathop{\mathrm{Coim}}}

% Boxtimes
%
\DeclareMathSymbol{\boxtimes}{\mathbin}{AMSa}{"02}

%
% Macros for moduli stacks/spaces
%
\def\QCohstack{\mathcal{QC}\!\mathit{oh}}
\def\Cohstack{\mathcal{C}\!\mathit{oh}}
\def\Spacesstack{\mathcal{S}\!\mathit{paces}}
\def\Quotfunctor{\mathrm{Quot}}
\def\Hilbfunctor{\mathrm{Hilb}}
\def\Curvesstack{\mathcal{C}\!\mathit{urves}}
\def\Polarizedstack{\mathcal{P}\!\mathit{olarized}}
\def\Complexesstack{\mathcal{C}\!\mathit{omplexes}}
% \Pic is the operator that assigns to X its picard group, usage \Pic(X)
% \Picardstack_{X/B} denotes the Picard stack of X over B
% \Picardfunctor_{X/B} denotes the Picard functor of X over B
\def\Pic{\mathop{\mathrm{Pic}}\nolimits}
\def\Picardstack{\mathcal{P}\!\mathit{ic}}
\def\Picardfunctor{\mathrm{Pic}}
\def\Deformationcategory{\mathcal{D}\!\mathit{ef}}
\usepackage{cleveref}
\crefname{originalequation}{original equation}{original equations}
\Crefname{originalequation}{Original equation}{Original equations}
\makeatletter
\newlabel{equation-j-p-shriek}{{0F4K}{source}{}{sourcepointer-1}{}}
\newlabel{equation-j-p-shriek@cref}{{[originalequation][0][]0F4K}{[1][0][]source}{}{}{}}
\makeatother
\crefname{originallemma}{original lemma}{original lemmas}
\Crefname{originallemma}{Original lemma}{Original lemmas}
\makeatletter
\newlabel{schemes-lemma-section-immersion}{{01KT}{source}{}{sourcepointer-2}{}}
\newlabel{schemes-lemma-section-immersion@cref}{{[originallemma][0][]01KT}{[1][0][]source}{}{}{}}
\makeatother
\title{A functorial injection from extension by zero into the direct image for a separated etale morphism}
\author{The Stacks Project Authors}
\date{stacks-project:etale-cohomology.tex:local-snapshot-20261011; GNU FDL 1.2 or later}
\begin{document}
\maketitle
\noindent\textit{Editorial note.} All delivered mathematical text is the author's own, copied from the pinned Stacks Project chapter file etale-cohomology.tex; no statement or proof step was rewritten, reordered or removed.\par\medskip
\noindent\textit{Editorial note.} The proof is the author's complete argument; the abelian case is written out and the module case is covered by the same construction.\par\medskip
\noindent\textit{Editorial note.} Every labelled reference occurring inside the counted spans is carried in assembly.external\_references as a non-load-bearing navigation pointer (same-chapter ones at their definition site, cross-chapter ones at the invocation inside the counted span), so no load-bearing mathematics is replaced by a stub. No bibliography citation occurs inside the counted spans.\par\medskip
\noindent\textit{Editorial note.} The checked chapter list for same-name collisions is the set of stacks chapters staged in this repository: chow, cohomology, divisors, etale-cohomology, flat, more-algebra, more-morphisms, morphisms, spaces-more-morphisms, varieties together with this card's chapter etale-cohomology.\par\medskip
\noindent\textit{Editorial note.} The preamble is the chapter's own native preamble excerpt already used for this source: the Xy-pic setup (native preamble lines 14-18) and the native theorem declarations with the mathematical macros (native preamble lines 154-217).\par\medskip
\noindent\textit{Allowed standard prerequisites (editorial).}\begin{itemize}
\item Extension by zero j\_! and direct image j\_* along a separated etale morphism. Same chapter, the functors j\_! and j\_* attached to j : U -> X and their description by sections and by compactly supported sections. Supplies the two functors and the map compared in the statement.
\item Separation of j and the description of sections of the direct image. Same chapter, the consequences of separatedness of j for the description of sections of j\_* and for the finiteness of the supports appearing in the colimits defining j\_!. Supplies the injectivity argument, in which a finite sum of sections with disjoint supports is examined.
\end{itemize}\medskip
\ifdefined\hypertarget\hypertarget{sourcepointer-1}{}\fi
\noindent\textit{Source reference.} Original equation 0F4K: sources/stacks-etale-cohomology/pinned-original-source.tex, source lines 11405--11405. Navigation pointer invoked at proof line 11530; the counted span contains the pointer only and the cited result is not reproduced inside it, so the pointer is registered as a non-load-bearing navigation stub in the established form for Stacks material. No bibliography citation is involved.\par\medskip
\ifdefined\hypertarget\hypertarget{sourcepointer-2}{}\fi
\noindent\textit{Source reference.} Original lemma 01KT: sources/stacks-etale-cohomology/pinned-original-source.tex, source lines 11545--11545. Navigation pointer invoked at proof line 11545; the counted span contains the pointer only and the cited result is not reproduced inside it, so the pointer is registered as a non-load-bearing navigation stub in the established form for Stacks material. No bibliography citation is involved.\par\medskip
\begin{lemma}
\label{lemma-shriek-into-star-separated-etale}
Let $j : U \to X$ be separated and \'etale. Then there is a functorial
injective map $j_!\mathcal{F} \to j_*\mathcal{F}$
on abelian sheaves and sheaves of $\Lambda$-modules.
\end{lemma}

\begin{proof}
We prove this in the case of abelian sheaves. Let us construct a canonical map
$$
j_{p!}\mathcal{F} \to j_*\mathcal{F}
$$
of abelian presheaves on $X_\etale$ for any abelian sheaf
$\mathcal{F}$ on $U_\etale$ where $j_{p!}$ is as in
(\ref{equation-j-p-shriek}). Sheafification of this map will
be the desired map $j_!\mathcal{F} \to j_*\mathcal{F}$.
Evaluating both sides on $V \to X$ \'etale we obtain
$$
j_{p!}\mathcal{F}(V) =
\bigoplus\nolimits_{\varphi : V \to U} \mathcal{F}(V \xrightarrow{\varphi} U)
\quad\text{and}\quad
j_*\mathcal{F}(V) = \mathcal{F}(V \times_X U)
$$
For each $\varphi$ we have an open and closed immersion
$$
\Gamma_\varphi = (1, \varphi) : V \longrightarrow V \times_X U
$$
over $U$. It is open as it is a morphism between schemes \'etale
over $U$ and it is closed as it is a section of a scheme separated
over $V$ (Schemes, Lemma \ref{schemes-lemma-section-immersion}).
Thus for a section $s_\varphi \in \mathcal{F}(V \xrightarrow{\varphi} U)$
there exists a unique section $s'_\varphi$ in $\mathcal{F}(V \times_X U)$
which pulls back to $s_\varphi$ by $\Gamma_\varphi$
and which restricts to zero on the complement of the image of $\Gamma_\varphi$.

\medskip\noindent
To show that our map is injective suppose that
$\sum_{i = 1, \ldots, n} s_{\varphi_i}$ is an element of
$j_{p!}\mathcal{F}(V)$ in the formula above maps to zero
in $j_*\mathcal{F}(V)$. Our task is to show that
$\sum_{i = 1, \ldots, n} s_{\varphi_i}$ restricts to zero
on the members of an \'etale covering of $V$.
Looking at all pairwise equalizers
(which are open and closed in $V$) of the morphisms
$\varphi_i : V \to U$ and working locally on $V$, we
may assume the images of the morphisms
$\Gamma_{\varphi_1}, \ldots, \Gamma_{\varphi_n}$
are pairwise disjoint. Since our assumption is that
$\sum_{i = 1, \ldots, n} s'_{\varphi_i} = 0$
we then immediately conclude that $s'_{\varphi_i} = 0$
for each $i$ (by the disjointness of the supports of these
sections), whence $s_{\varphi_i} = 0$ for all $i$ as desired.
\end{proof}

\end{document}
